\begin{proofbox}
\textbf{统一记号与前提。}

设 $F_1=(-c,0),\ F_2=(c,0)$，
$P=(x_0,y_0)\in H$ 且 $y_0\ne0$，
$M=(m,0)$ 是 $\angle F_1PF_2$ 内角平分线的轴上交点。
令
\[
r_1=PF_1,\qquad r_2=PF_2,\qquad c^2=a^2+b^2.
\]
由于 $P$ 非顶点，$x_0>a$ 或 $x_0<-a$，而 $M$ 位于开线段 $F_1F_2$ 上。
以下给出\textbf{五种方法}，排序考虑入门负担与考场速度。

\medskip
\textbf{方法一：最低门槛【角平分线定理＋平方差；首推独立证明】。}

\textbf{第1步：角平分线定理。}
\[
\frac{c+m}{c-m}=\frac{r_1}{r_2}
\quad\Longrightarrow\quad
r_2(c+m)=r_1(c-m),
\]
于是
\[
\boxed{m=c\frac{r_1-r_2}{r_1+r_2}}.\tag{1}
\]
\textbf{第2步：焦距平方相减。}
\[
\begin{aligned}
r_1^2-r_2^2
&=(x_0+c)^2+y_0^2-(x_0-c)^2-y_0^2\\
&=4cx_0.
\end{aligned}
\]
即
\[
\boxed{x_0=\frac{(r_1-r_2)(r_1+r_2)}{4c}}.\tag{2}
\]
\textbf{第3步：乘起来，使用双曲线第一定义。}
\[
\begin{aligned}
mx_0
&=\frac{c(r_1-r_2)}{r_1+r_2}
  \cdot\frac{(r_1-r_2)(r_1+r_2)}{4c}\\
&=\frac{(r_1-r_2)^2}{4}
=\frac{(2a)^2}{4}=\boxed{a^2}.
\end{aligned}
\]
\textbf{优势：}完全不必区分左右支；不背焦半径；只用角平分线定理、
两点距离与平方差公式。\textbf{关键切入：}两个焦距是“和、差”，而 $x_0$ 是“平方差”。

\medskip
\textbf{方法二：最快代数【角平分线定理＋焦半径】。}

先沿用式 (1)。令 $e=c/a>1$，记
$\sigma=\operatorname{sgn}(x_0)\in\{1,-1\}$。
双曲线左右两支的焦半径可以统一写成
\[
\boxed{r_1=e|x_0|+\sigma a,\qquad
r_2=e|x_0|-\sigma a}.
\]
核对：右支 $\sigma=1$ 时 $r_1=ex_0+a,\ r_2=ex_0-a$；
左支 $\sigma=-1$ 时 $r_1=-ex_0-a,\ r_2=-ex_0+a$。
因此
\[
\begin{aligned}
m
&=c\frac{2\sigma a}{2e|x_0|}
=\frac{ca\sigma}{e|x_0|}
=\frac{a^2\sigma}{|x_0|}
=\boxed{\frac{a^2}{x_0}}.
\end{aligned}
\]
\textbf{优势：}已熟悉焦半径时只需代入一次。
\textbf{代价：}若左右焦点匹配或左右支的符号记错，很容易出错。

\medskip
\textbf{方法三：最快识别【已知切线平分焦点三角形的内角】。}

若已掌握双曲线的\textbf{切线--焦点角平分线性质}，
$PM$ 就是 $H$ 在 $P$ 处的切线。
切线方程为
\[
\frac{x_0x}{a^2}-\frac{y_0y}{b^2}=1.
\]
令 $y=0$，得到
\[
\frac{x_0m}{a^2}=1
\quad\Longrightarrow\quad
\boxed{m=\frac{a^2}{x_0}}.
\]
\textbf{优势：}看出切线后，求值几乎一步完成。
\textbf{注意：}本方法\emph{调用了已知的切线--角平分线性质}，
不能再把这一步当作对该性质的独立证明；若还没掌握该性质，用方法一建立它。

\medskip
\textbf{方法四：几何深化【内切圆切线长＋角平分线定理】。}

作 $\triangle PF_1F_2$ 的内切圆，切 $F_1F_2$ 于 $T(t,0)$，
切 $PF_1$、$PF_2$ 于 $U,V$。令 $PU=PV=\lambda>0$。
同一点的两条切线段等长，所以
\[
r_1=F_1T+\lambda,\qquad
r_2=F_2T+\lambda.
\]
故
\[
F_1T-F_2T=r_1-r_2=2\sigma a,
\qquad F_1T+F_2T=2c,
\]
\[
F_1T=c+\sigma a,\qquad
F_2T=c-\sigma a,
\qquad\boxed{T=(\sigma a,0)}.
\]
这先证明了\textbf{内切圆与焦点轴的切点始终为对应分支顶点}。
并且
\[
r_1=\lambda+c+\sigma a,\qquad
r_2=\lambda+c-\sigma a,
\]
由距离平方差
\[
4cx_0=r_1^2-r_2^2
=(2\sigma a)\cdot2(\lambda+c),
\]
得
\[
\lambda+c=\frac{c|x_0|}{a}.
\]
最后用角平分线定理的式 (1)：
\[
m=c\frac{r_1-r_2}{r_1+r_2}
=c\frac{2\sigma a}{2(\lambda+c)}
=\boxed{\frac{a^2}{x_0}}.
\]
\textbf{优势：}揭示定切点 $T=\pm a$，几何画面非常有记忆点。
\textbf{代价：}为了进一步求 $M$，还需连接焦距平方差或其他长度关系；
不如方法一直接。切记 $T$ 与 $M$ 通常不重合。

\medskip
\textbf{方法五：半角三角函数【单位方向向量的内角平分线】。}

设两条从 $P$ 指向焦点的\textbf{单位方向向量}为
\[
\boldsymbol{u}_1=\frac{(-c-x_0,-y_0)}{r_1},
\qquad
\boldsymbol{u}_2=\frac{(c-x_0,-y_0)}{r_2}.
\]
若顶角 $\gamma=\angle F_1PF_2\in(0,\pi)$，
两个单位向量相加得到
\[
\boldsymbol{u}_1+\boldsymbol{u}_2
=2\cos\frac{\gamma}{2}\,\boldsymbol{w},
\]
其中 $\boldsymbol{w}$ 是\textbf{内角平分线的单位方向}。
这一等式来自\emph{半角余弦公式}；因此
$\boldsymbol{v}=\boldsymbol{u}_1+\boldsymbol{u}_2$ 就沿 $PM$。

设 $M=P+k\boldsymbol v$。利用 $M$ 的纵坐标为 $0$，
\[
0=y_0-k y_0\Bigl(\frac1{r_1}+\frac1{r_2}\Bigr)
\quad\Longrightarrow\quad
k=\frac{r_1r_2}{r_1+r_2}.
\]
代入横坐标：
\[
\begin{aligned}
m
&=x_0+k\Bigl(-\frac{x_0+c}{r_1}+\frac{c-x_0}{r_2}\Bigr)\\
&=x_0+\frac{-(x_0+c)r_2+(c-x_0)r_1}{r_1+r_2}\\
&=c\frac{r_1-r_2}{r_1+r_2}.
\end{aligned}
\]
再用 $r_1^2-r_2^2=4cx_0$ 与 $|r_1-r_2|=2a$，
便得 $mx_0=a^2$。
\textbf{优势：}直接由\emph{半角方向}建立直线，无需预先调用角平分线定理。
\textbf{代价：}向量/三角运算更多，不适合追求最低认知负担的首解。

\medskip
\textbf{补充验算：从 $x_0m=a^2$ 独立推出切线方程。}

因 $y_0\ne0$，经过 $P(x_0,y_0)$、$M(a^2/x_0,0)$ 的直线满足
\[
y(x_0^2-a^2)=y_0(xx_0-a^2).
\]
又因 $P\in H$，
\[
\frac{x_0^2}{a^2}-\frac{y_0^2}{b^2}=1
\quad\Longrightarrow\quad
x_0^2-a^2=\frac{a^2y_0^2}{b^2}.
\]
代回并除以 $a^2y_0$：
\[
\frac{y_0y}{b^2}=\frac{x_0x}{a^2}-1,
\]
即
\[
\boxed{PM:\frac{x_0x}{a^2}-\frac{y_0y}{b^2}=1}.
\]
这是双曲线在 $P$ 处的标准切线方程，
因而\textbf{用方法一不仅证明截距，还独立证明了内角平分线即切线}。

\medskip
\textbf{选法建议：}需要\emph{最低门槛独立证明}用方法一；
已熟悉焦半径且追求代数速度用方法二；
切线性质已是已知结论时用方法三；
要记住内切圆固定切点用方法四；
研究内角平分线的半角方向用方法五。
\end{proofbox}
